% Options for packages loaded elsewhere
\PassOptionsToPackage{unicode}{hyperref}
\PassOptionsToPackage{hyphens}{url}
\PassOptionsToPackage{dvipsnames,svgnames,x11names}{xcolor}
%
\documentclass[
  13pt,
  a4paperpaper,
  12pt]{article}
\usepackage{amsmath,amssymb}
\usepackage{setspace}
\usepackage{iftex}
\ifPDFTeX
  \usepackage[T1]{fontenc}
  \usepackage[utf8]{inputenc}
  \usepackage{textcomp} % provide euro and other symbols
\else % if luatex or xetex
  \usepackage{unicode-math} % this also loads fontspec
  \defaultfontfeatures{Scale=MatchLowercase}
  \defaultfontfeatures[\rmfamily]{Ligatures=TeX,Scale=1}
\fi
\usepackage{lmodern}
\ifPDFTeX\else
  % xetex/luatex font selection
  \setmainfont[]{Times New Roman}
  \setmonofont[]{Courier New}
\fi
% Use upquote if available, for straight quotes in verbatim environments
\IfFileExists{upquote.sty}{\usepackage{upquote}}{}
\IfFileExists{microtype.sty}{% use microtype if available
  \usepackage[]{microtype}
  \UseMicrotypeSet[protrusion]{basicmath} % disable protrusion for tt fonts
}{}
\makeatletter
\@ifundefined{KOMAClassName}{% if non-KOMA class
  \IfFileExists{parskip.sty}{%
    \usepackage{parskip}
  }{% else
    \setlength{\parindent}{0pt}
    \setlength{\parskip}{6pt plus 2pt minus 1pt}}
}{% if KOMA class
  \KOMAoptions{parskip=half}}
\makeatother
\usepackage{xcolor}
\usepackage[margin=1.5cm,top=1.5cm,bottom=1.5cm,left=1.5cm,right=1.5cm,includeheadfoot]{geometry}
\usepackage{listings}
\newcommand{\passthrough}[1]{#1}
\lstset{defaultdialect=[5.3]Lua}
\lstset{defaultdialect=[x86masm]Assembler}
\setlength{\emergencystretch}{3em} % prevent overfull lines
\providecommand{\tightlist}{%
  \setlength{\itemsep}{0pt}\setlength{\parskip}{0pt}}
\setcounter{secnumdepth}{3}
% Minimal LaTeX preamble fragment for Pandoc-built documents
% Avoid \documentclass or loading hyperref/cleveref here (they are managed by the build pipeline)
% Provide safe unicode declarations and small custom macros only
\DeclareUnicodeCharacter{03BC}{\(\mu\)}
\DeclareUnicodeCharacter{03BB}{\(\lambda\)}
\DeclareUnicodeCharacter{03BD}{\(\nu\)}
\DeclareUnicodeCharacter{03C0}{\(\pi\)}
\DeclareUnicodeCharacter{03B5}{\(\epsilon\)}
\DeclareUnicodeCharacter{03B4}{\(\delta\)}

% Map common Unicode subscripts and degree symbol to LaTeX-friendly macros
\DeclareUnicodeCharacter{2080}{\textsubscript{0}}
\DeclareUnicodeCharacter{2081}{\textsubscript{1}}
\DeclareUnicodeCharacter{2082}{\textsubscript{2}}
\DeclareUnicodeCharacter{2083}{\textsubscript{3}}
\DeclareUnicodeCharacter{2084}{\textsubscript{4}}
\DeclareUnicodeCharacter{2085}{\textsubscript{5}}
\DeclareUnicodeCharacter{2086}{\textsubscript{6}}
\DeclareUnicodeCharacter{2087}{\textsubscript{7}}
\DeclareUnicodeCharacter{2088}{\textsubscript{8}}
\DeclareUnicodeCharacter{2089}{\textsubscript{9}}
\DeclareUnicodeCharacter{00B0}{\textdegree}

% Counter behaviour
\usepackage{chngcntr}
\counterwithout{figure}{section}
\counterwithout{table}{section}
\counterwithout{equation}{section}

% Caption formatting
\usepackage[font=small,labelfont=bf,textfont=it]{caption}
\captionsetup{justification=raggedright,singlelinecheck=false}
\captionsetup[figure]{position=bottom,skip=10pt}
\captionsetup[table]{position=top,skip=10pt}
\AtBeginEnvironment{figure}{\raggedright}

% Titlesec formatting
\usepackage{titlesec}
\titlespacing{\section}{0pt}{12pt}{6pt}
\titlespacing{\subsection}{0pt}{10pt}{4pt}
\titlespacing{\subsubsection}{0pt}{8pt}{3pt}
\titleformat{\section}{\large\bfseries}{\thesection}{1em}{}
\titleformat{\subsection}{\normalsize\bfseries}{\thesubsection}{1em}{}
\titleformat{\subsubsection}{\normalsize\itshape}{\thesubsubsection}{1em}{}

% Ensure each top-level section starts on a new page
\let\oldsection\section
\renewcommand{\section}{\clearpage\oldsection}

% Graphics support for figures
\usepackage{graphicx}

% Small visual tweaks
\usepackage{xcolor}
\definecolor{codebg}{RGB}{245,245,245}
\definecolor{codeborder}{RGB}{200,200,200}
\definecolor{codefg}{RGB}{50,50,50}

% Setup listings safely (don't override global styles)
\usepackage{listings}
\lstset{backgroundcolor=\color{codebg},basicstyle=\small\ttfamily\color{codefg},frame=single,framerule=0.5pt,framesep=5pt}

% Keep title metadata to be inserted by build system
\ifLuaTeX
  \usepackage{selnolig}  % disable illegal ligatures
\fi
\IfFileExists{bookmark.sty}{\usepackage{bookmark}}{\usepackage{hyperref}}
\IfFileExists{xurl.sty}{\usepackage{xurl}}{} % add URL line breaks if available
\urlstyle{same}
\hypersetup{
  colorlinks=true,
  linkcolor={blue},
  filecolor={blue},
  citecolor={blue},
  urlcolor={blue},
  pdfcreator={LaTeX via pandoc}}

\author{}
\date{}

\begin{document}

{
\hypersetup{linkcolor=black}
\setcounter{tocdepth}{3}
\tableofcontents
}
\setstretch{1.0}
\section{Conclusion}\label{sec:conclusion}

\subsection{Summary of findings}\label{summary-of-findings}

We present a reproducible computational framework that implements,
tests, and evaluates the vibrational hypothesis for insect olfaction.
Integrating morphology, spectroscopy, neural timing, and environmental
modeling, the framework produces quantitative predictions and explicit
falsifiers suitable for experimental validation.

\subsubsection{Key innovation}\label{key-innovation}

All predictions are anchored in deterministic, unit‑tested code with
documented case studies and reproducible figure generation. This ensures
traceability from equations to figures and tests.

\subsubsection{Empirical highlights}\label{empirical-highlights}

\begin{enumerate}
\def\labelenumi{\arabic{enumi}.}
\tightlist
\item
  Morphology: Sensilla dimensions (typically 2--160 \(\mu\mathrm{m}\)
  across sensillum types) frequently align with predicted resonant
  wavelengths (example correlations r \(\approx 0.85\) across sampled
  taxa); computations reproduced by
  \passthrough{\lstinline!src/sensilla.py::analyze\_sensilla\_dimensions!}.
\item
  Neural timing: Meta‑analysis and modeling show ORN latencies (1--5 ms;
  minimum \textasciitilde3 ms reported) that can be reconciled with an
  early IR‑detection stage (see
  \passthrough{\lstinline!src/core.py::calculate\_response\_time\_improvement!}).
\item
  Behavior: Specialized IR sensors and CHC‑dependent behaviors are
  consistent with frequency‑specific detection thresholds reported in
  the literature.
\item
  Spectroscopy: Automated CHC peak detection identifies bands within
  atmospheric windows with (\pm)0.1 \(\mu\mathrm{m}\) resolution
  (implemented in \passthrough{\lstinline!src/spectroscopy.py!}).
\end{enumerate}

\subsection{Future Research Directions and
Applications}\label{future-research-directions-and-applications}

\subsubsection{Immediate Experimental
Priorities}\label{immediate-experimental-priorities}

\textbf{High-Impact Validation Studies:}

\begin{itemize}
\tightlist
\item
  \textbf{Single-sensillum electrophysiology} with quantum cascade laser
  stimulation (2--25 \(\mu\mathrm{m}\))
\item
  \textbf{Behavioral IR-only orientation assays} using narrowband LEDs
  with thermal controls
\item
  \textbf{Cross-taxa morphometric surveys} with SEM and resonance
  correlation analysis
\item
  \textbf{High-temporal-resolution neural recordings} (sub-ms) to
  resolve detection components
\end{itemize}

\textbf{Advanced Instrumentation Development:}

\begin{itemize}
\tightlist
\item
  \textbf{Tunable IR sources} with sub-micron wavelength precision
\item
  \textbf{Thermal imaging integration} with neural recordings
\item
  \textbf{Multi-scale sensing platforms} combining electromagnetic and
  molecular detection
\end{itemize}

\subsubsection{Computational and Theoretical
Extensions}\label{computational-and-theoretical-extensions}

\textbf{Enhanced Modeling Frameworks:}

\begin{itemize}
\tightlist
\item
  \textbf{3D electromagnetic simulations} of complete antenna systems
\item
  \textbf{Machine learning classifiers} for spectral feature recognition
\item
  \textbf{Climate model integration} for environmental robustness
  predictions
\item
  \textbf{Quantum mechanical extensions} incorporating coherence and
  entanglement effects
\end{itemize}

\textbf{Information Processing Analysis:}

\begin{itemize}
\tightlist
\item
  \textbf{Channel capacity optimization} under realistic noise
  conditions
\item
  \textbf{Population coding strategies} for IR signal processing
\item
  \textbf{Neural network models} integrating electromagnetic and
  molecular pathways
\end{itemize}

\subsubsection{Technological Translation and
Applications}\label{technological-translation-and-applications}

\textbf{Biomimetic Sensor Development:}

\begin{itemize}
\tightlist
\item
  \textbf{IR detection arrays} inspired by insect sensilla geometry
\item
  \textbf{Wearable environmental monitors} for chemical detection
\item
  \textbf{Autonomous navigation systems} using vibrational cues
\item
  \textbf{Medical diagnostic tools} based on molecular vibration sensing
\end{itemize}

\textbf{Agricultural and Environmental Applications:}

\begin{itemize}
\tightlist
\item
  \textbf{Precision pest monitoring} using species-specific IR
  signatures
\item
  \textbf{Smart trap systems} with adaptive wavelength selection
\item
  \textbf{Pollination monitoring} through floral scent IR profiling
\item
  \textbf{Conservation tools} for endangered species detection
\end{itemize}

\subsubsection{Fundamental Scientific
Implications}\label{fundamental-scientific-implications}

\textbf{Sensory Biology Advancements:}

\begin{itemize}
\tightlist
\item
  \textbf{Multi-modal integration} models combining electromagnetic and
  molecular sensing
\item
  \textbf{Evolutionary convergence} studies across taxa with IR
  detection
\item
  \textbf{Quantum effects} in biological signal processing
\item
  \textbf{Sensory ecology} incorporating infrared communication channels
\end{itemize}

\textbf{Broader Theoretical Implications:}

\begin{itemize}
\tightlist
\item
  \textbf{Unified theory} of olfaction integrating shape and vibration
\item
  \textbf{Electromagnetic sensing} in other biological systems
\item
  \textbf{Quantum biology} extensions to sensory processing
\item
  \textbf{Bio-inspired engineering} principles for next-generation
  sensors
\end{itemize}

\subsection{Implementation Roadmap}\label{implementation-roadmap}

\subsubsection{Phase 1: Core Validation (6--12
months)}\label{phase-1-core-validation-612-months}

\begin{itemize}
\tightlist
\item
  Complete single-sensillum IR sensitivity studies
\item
  Establish behavioral IR-only assay protocols
\item
  Validate morphometric correlations across 10+ species
\item
  Develop standardized thermal control methodologies
\end{itemize}

\subsubsection{Phase 2: Technology Development (1--2
years)}\label{phase-2-technology-development-12-years}

\begin{itemize}
\tightlist
\item
  Prototype biomimetic IR sensor arrays
\item
  Integrate multi-modal sensing platforms
\item
  Develop field-deployable instrumentation
\item
  Create comprehensive spectral databases
\end{itemize}

\subsubsection{Phase 3: Applied Translation (2--3
years)}\label{phase-3-applied-translation-23-years}

\begin{itemize}
\tightlist
\item
  Commercial sensor platform development
\item
  Agricultural pest management systems
\item
  Environmental monitoring networks
\item
  Medical diagnostic applications
\end{itemize}

\subsection{Reproducibility and Knowledge
Transfer}\label{reproducibility-and-knowledge-transfer}

The Appendices and \passthrough{\lstinline!src/!} modules provide
complete computational anchors for all claims, enabling:

\begin{itemize}
\tightlist
\item
  \textbf{Independent validation} of theoretical predictions
\item
  \textbf{Experimental protocol reproduction} with detailed
  specifications
\item
  \textbf{Technology transfer} to industrial and academic partners
\item
  \textbf{Educational resources} for training next-generation
  researchers
\end{itemize}

This integrated framework establishes a foundation for both fundamental
understanding of insect sensory mechanisms and practical applications in
biomimetic sensing technology.

\end{document}
